%
% display_crash_loop_earlyexit_buggy.tex
%
% FIDELITY NEGATIVE CONTROL #2 for display_crash_loop.tex: the STABLE-INTERVAL
% EXIT clause struck.  This is the canonical corrected model with one change --
% IsolExit's guard drops "crasher <: quarantined", so the worker may leave
% isolation EARLY (the "one clean pass" exit) and wipe the tallies of scenes
% still in play.  An intermittent crasher then never accumulates: it crashes
% the batch (tallied), enters isolation, renders cleanly on its probe, IsolExit
% fires early and wipes its tally, and the next batched death restarts the
% cycle.  crashes grows without bound; ProB reaches
% "crashes > 1 + (THRESHOLD-1)*card(crasher)".  This is the intermittent-crasher
% loop the STABLE_INTERVAL >= ATTRIBUTION_WINDOW tie prevents.
%
% Do NOT treat this as a design artifact; the canonical, verified model is
% display_crash_loop.tex.  See also display_crash_loop_buggy.tex, the control
% for the quarantine clause.
%
% Type-check:  fuzz -t docs/display_crash_loop_earlyexit_buggy.tex
% Model-check (must report: goal FOUND -- the loop):
%   probcli docs/display_crash_loop_earlyexit_buggy.tex -model_check -nodead \
%       -goal "crashes>1+(THRESHOLD-1)*card(crasher)" \
%       -p DEFAULT_SETSIZE 2 -p MAX_INITIALISATIONS 20 -p MAXINT 8
%

\documentclass[a4paper,10pt,fleqn]{article}
\usepackage[margin=1in]{geometry}
\usepackage{fuzz}
\usepackage[colorlinks=true,linkcolor=blue,citecolor=blue,urlcolor=blue]{hyperref}

\begin{document}

\title{Crash-Loop Quarantine: Fidelity Control \#2 (early isolation exit)}
\author{Buggy variant --- isolation left before the crashers are caught}
\date{July 2026}
\maketitle

% ============================================================
\section{Introduction}
% ============================================================

A scene that passes Hub-side self-validation can still crash the Display's
renderer.  When it does, the Display process dies --- but the scene does not
die with it.  It stays live in \texttt{HubDisplay}, the authoritative store, so
the Hub's crash recovery re-marks it for repaint, the respawned Display receives
it, and it crashes again.  Each respawn opens a fresh window that steals
keyboard focus.  Nothing in this cycle attributes the death to the scene that
caused it, so nothing stops the loop; the operator observed it live during the
B5 demo, once every ten seconds, until the scenes were cleared by hand.

The fix is a \emph{quarantine} rule.  Every Display death is attributed to the
scenes that were in flight when it happened: in batching mode that is the whole
coalesced batch, in isolation mode a single probed scene.  A scene quarantined
after \texttt{THRESHOLD} attributed deaths stops being replicated while staying
in the store for inspection.  Two subtleties make the rule sound against an
\emph{intermittently} poisonous scene --- one that crashes the renderer on only
some renders.  First, no death is free: even a batched death advances the tally,
so a scene that crashes only while coalesced cannot escape.  Second, isolation is
left only after a stable no-death interval, not after one clean pass, so a flaky
scene is pursued across its clean renders until it is caught, and an innocent
scene's lone batched increment ages out before a second one can land.

This spec models that lifecycle as a finite state machine, so ProB can confirm
the three properties that matter --- a quarantined scene is never replicated, the
number of crashes is bounded (the loop terminates, even for a flaky scene), and
an innocent scene is never quarantined --- and can reproduce the loop when either
corrective clause is removed.

This spec is a companion to \texttt{display\_lifecycle.tex}, not an extension of
it.  That spec proves that exactly one Display process owns the socket path under
a two-lock discipline; this spec takes that single-owner Display as given and
models what happens when the scene it renders is poisonous.  The two concerns are
orthogonal: the bind-race spec governs \emph{who owns the socket}, this spec
governs \emph{whether a poison scene can loop the Display}.

% ============================================================
\section{Basic Types}
% ============================================================

\begin{zed}
  [SCENE]
\end{zed}

\texttt{SCENE} is the set of scenes an agent has installed in the store.  Two
scenes suffice to exhibit a poison scene alongside an innocent one; three
exercise attribution against several crashers at once.  The carrier is bounded by
ProB's \texttt{DEFAULT\_SETSIZE}.

% ============================================================
\section{Free Types}
% ============================================================

\begin{zed}
  DISPLAY ::= up | down
\end{zed}

The state of the single Display process: \texttt{up} while serving,
\texttt{down} after a renderer crash and before respawn.  That there is exactly
one Display process is proved in \texttt{display\_lifecycle.tex}; this model
represents that one process's liveness.

\begin{zed}
  MODE ::= batching | isolating
\end{zed}

The replicator's send mode.  \texttt{batching}: dirty scenes are coalesced and
sent together.  \texttt{isolating}: each scene is sent in its own send, so a
death has a single suspect.  The worker starts \texttt{batching}, switches to
\texttt{isolating} on the first death, and returns to \texttt{batching} only
after a stable no-death interval.

% ============================================================
\section{Constants}
% ============================================================

\begin{axdef}
  THRESHOLD : \nat_1
\where
  THRESHOLD = 2
\end{axdef}

\texttt{THRESHOLD} is the number of attributed deaths a scene must cause before
it is quarantined.  It is two, not one: a single death admits a non-scene cause
(memory pressure, a driver fault) that coincided with a render, and must not
quarantine a scene.  Requiring two excludes the one-off, while a genuinely poison
scene reaches it deterministically.

% ============================================================
\section{State}
% ============================================================

The state is flat: one schema, six components.  Its predicate carries only the
\emph{coherence} constraint that a quarantined scene has reached the threshold ---
true in every operation by construction.  The safety properties
(Section~\ref{sec:inv}) are deliberately \emph{absent} here, so a violating state
stays reachable and ProB can find it.  In particular, \texttt{quarantined
$\subseteq$ crasher} --- innocent scenes are never quarantined --- is \emph{not}
asserted here; it is a property to be model-checked, not a constraint to be
assumed.

\begin{schema}{Sys}
  disp : DISPLAY \\
  mode : MODE \\
  crasher : \power SCENE \\
  quarantined : \power SCENE \\
  tally : SCENE \fun \nat \\
  crashes : \nat
\where
  \forall s : SCENE @ s \in quarantined \implies tally~s \geq THRESHOLD
\end{schema}

\texttt{crasher} is the set of scenes whose render \emph{can} crash the Display
--- an environmental fact fixed for the run.  A deterministic crasher crashes on
every render; an intermittent one crashes on only some, which the model captures
by letting a crasher's isolation probe either crash or render cleanly (operations
\texttt{IsolCrash} and \texttt{IsolClean}).  \texttt{quarantined} is the set the
Hub has stopped replicating, \texttt{tally} counts attributed deaths per scene,
and \texttt{crashes} counts Display deaths in total --- the quantity whose
boundedness is the no-infinite-respawn property.

% ============================================================
\section{Initialization}
% ============================================================

\begin{schema}{Init}
  Sys'
\where
  disp' = up \\
  mode' = batching \\
  crasher' \subseteq SCENE \\
  quarantined' = \emptyset \\
  tally' = (\lambda s : SCENE @ 0) \\
  crashes' = 0
\end{schema}

The Display starts up, the worker batching, nothing quarantined, every tally
zero.  \texttt{crasher} is left free: ProB explores every valuation --- no
crasher, all crashers, and every mix, including an innocent scene beside a
crasher --- in one run.

% ============================================================
\section{Operations}
% ============================================================

In batching mode a coalesced send crashes the Display when any live,
non-quarantined crasher is in the batch.  The death is attributed to the whole
batch --- every non-quarantined scene has its tally incremented --- and the
worker switches to isolation.  Attributing the batched death is what stops a
scene that crashes only while coalesced from escaping the tally.

\begin{schema}{BatchCrash}
  \Delta Sys \\
  c? : SCENE
\where
  disp = up \\
  mode = batching \\
  c? \in crasher \\
  c? \notin quarantined \\
  disp' = down \\
  crashes' = crashes + 1 \\
  mode' = isolating \\
  tally' = \\
  \quad~ (\lambda s : SCENE @ \\
  \quad\quad~ \IF s \notin quarantined \THEN tally~s + 1 \ELSE tally~s) \\
  quarantined' = \\
  \quad~ quarantined \cup \\
  \quad\quad~ \{ s : SCENE | s \notin quarantined \land tally~s + 1 \geq THRESHOLD \} \\
  crasher' = crasher
\end{schema}

In isolation mode each scene is probed alone.  An intermittent crasher may crash
this probe or render cleanly.  A crash is attributed to that one scene and
quarantines it at the threshold; the worker stays in isolation.

\begin{schema}{IsolCrash}
  \Delta Sys \\
  s? : SCENE
\where
  disp = up \\
  mode = isolating \\
  s? \in crasher \\
  s? \notin quarantined \\
  disp' = down \\
  crashes' = crashes + 1 \\
  tally' = tally \oplus \{ s? \mapsto tally~s? + 1 \} \\
  quarantined' = \\
  \quad~ \IF tally~s? + 1 \geq THRESHOLD \\
  \quad~ \THEN quarantined \cup \{ s? \} \\
  \quad~ \ELSE quarantined \\
  mode' = mode \\
  crasher' = crasher
\end{schema}

The guard \texttt{s? $\notin$ quarantined} is the quarantine defence: a scene
already quarantined can never be the subject of \texttt{IsolCrash}, so it never
crashes the Display again.  A clean isolation probe of an intermittent crasher
changes nothing and, crucially, does \emph{not} leave isolation:

\begin{schema}{IsolClean}
  \Xi Sys \\
  s? : SCENE
\where
  disp = up \\
  mode = isolating \\
  s? \in crasher \\
  s? \notin quarantined
\end{schema}

The worker leaves isolation only when the Display has served a stable interval
with no death.  In the model that interval elapses exactly when no
non-quarantined crasher remains --- there is nothing left to crash --- so the
guard is \texttt{crasher $\subseteq$ quarantined}.  On exit the tallies of scenes
still in play are cleared, modelling the attribution window ageing out every
prior death once the display has been death-free for at least a window:

\begin{schema}{IsolExit}
  \Delta Sys
\where
  disp = up \\
  mode = isolating \\
  mode' = batching \\
  tally' = \\
  \quad~ (\lambda s : SCENE @ \IF s \in quarantined \THEN tally~s \ELSE 0) \\
  disp' = disp \\
  crasher' = crasher \\
  quarantined' = quarantined \\
  crashes' = crashes
\end{schema}

A crashed Display is respawned.  This corresponds to the reap-then-ensure pair
that \texttt{display\_lifecycle.tex} proves yields exactly one socket owner:

\begin{schema}{Respawn}
  \Delta Sys
\where
  disp = down \\
  disp' = up \\
  mode' = mode \\
  crasher' = crasher \\
  quarantined' = quarantined \\
  tally' = tally \\
  crashes' = crashes
\end{schema}

The respawn backoff of the design --- which paces successive respawns so the
pre-quarantine deaths are not a rapid focus-stealing burst --- is a timing
policy, not a state-safety property: it changes \emph{when} \texttt{Respawn}
fires, never \emph{whether} a reachable state violates an invariant.  It is
specified in the design document, not modelled.

% ============================================================
\section{Invariants}
\label{sec:inv}
% ============================================================

Four properties must hold across all reachable states.  None is placed in the
\texttt{Sys} schema: a predicate placed there would be enforced as a guard on
every successor, silently disabling any operation that would break it and thereby
masking the very defect this model exists to catch.

\paragraph{Invariant 1 --- quarantined is never replicated.}
No quarantined scene is ever rendered.  Established by construction:
\texttt{BatchCrash}, \texttt{IsolCrash}, and \texttt{IsolClean} --- the only
operations that send a scene to the Display --- are each guarded by
\texttt{$\notin$ quarantined} on the scene they touch.  ProB's operation-coverage
report corroborates that none fires against a quarantined scene.

\paragraph{Invariant 2 --- crashes are bounded (no infinite respawn).}
The total number of Display crashes never exceeds one batched death plus the
isolated deaths needed to drive each crasher to the threshold:
\[
  crashes \leq 1 + (THRESHOLD - 1) * \# crasher.
\]
Exactly one \texttt{BatchCrash} occurs per run: it fires only in batching mode,
and it switches the worker to isolation, which is left (via \texttt{IsolExit})
only once every crasher is quarantined --- after which no non-quarantined crasher
remains to trigger a second \texttt{BatchCrash}.  That one batched death advances
every crasher's tally to one, so each crasher then needs at most
\texttt{THRESHOLD} $-$ \texttt{1} \texttt{IsolCrash} deaths to reach the
threshold.  \texttt{IsolClean} adds no crash, so an \emph{intermittent} crasher
--- which interleaves clean probes --- reaches the same bound: isolation persists
across its clean renders, so its crashes still accumulate.  This is the
no-infinite-respawn property.

\paragraph{Invariant 3 --- an innocent scene is never quarantined.}
A scene that is not a crasher is never quarantined:
\[
  quarantined \subseteq crasher.
\]
An innocent scene gains a tally only from the single batched death (it is never a
crashing isolation singleton), so its tally reaches at most one, below the
threshold of two.  Because \texttt{quarantined $\subseteq$ crasher} is deliberately
kept out of the \texttt{Sys} predicate, this is a genuine reachability check, not
an assumed constraint: \texttt{BatchCrash} \emph{could} quarantine any scene that
reached the threshold, and ProB confirms an innocent one never does.

\paragraph{Invariant 4 --- single Display.}
At most one Display process is live.  Structural here (\texttt{disp} is a single
\texttt{DISPLAY}-valued component); its real content --- that reap/ensure yields
exactly one socket owner --- is proved in \texttt{display\_lifecycle.tex} and
inherited by composition.  \texttt{Respawn} abstracts that verified pair.

% ============================================================
\section{Verification}
\label{sec:verify}
% ============================================================

Type-checking is a \texttt{fuzz} pass:
\begin{verbatim}
  fuzz -t docs/display_crash_loop.tex
\end{verbatim}

Invariants 2 and 3 are state properties, checked by asking ProB whether their
negation is \emph{reachable}.  A goal that is not found means the invariant holds
on every reachable state:
\begin{verbatim}
  # Invariant 2 (must report: goal NOT found)
  probcli docs/display_crash_loop.tex -model_check -nodead \
    -goal "crashes>1+(THRESHOLD-1)*card(crasher)" \
    -p DEFAULT_SETSIZE 2 -p MAX_INITIALISATIONS 20

  # Invariant 3 (must report: goal NOT found)
  probcli docs/display_crash_loop.tex -model_check -nodead \
    -goal "quarantined - crasher /= {}" \
    -p DEFAULT_SETSIZE 2 -p MAX_INITIALISATIONS 20
\end{verbatim}

The model also reaches the stable fixed point in which every crasher is
quarantined, the worker is back in batching mode, and the Display is up --- the
state the loop settles into.  That this state is reachable is a positive goal ProB
finds:
\begin{verbatim}
  # Stable state reachable (must report: goal FOUND)
  probcli docs/display_crash_loop.tex -model_check -nodead \
    -goal "disp=up & mode=batching & crasher/={} & crasher<:quarantined" \
    -p DEFAULT_SETSIZE 2 -p MAX_INITIALISATIONS 20
\end{verbatim}

Invariant 1 is a transition property, established by the guards of
\texttt{BatchCrash}, \texttt{IsolCrash}, and \texttt{IsolClean} and corroborated
by ProB's operation-coverage report.  Invariant 4 is inherited from
\texttt{display\_lifecycle.tex}.

The \texttt{-nodead} flag is used deliberately.  Unlike the bind-race spec, whose
liveness depends on never deadlocking, this model is \emph{meant} to terminate:
the state in which every crasher is quarantined and the Display is up has no
enabled operation, and that quiescence is the desired outcome, not a fault.  The
\texttt{MAX\_INITIALISATIONS} bound is raised past the default four because
\texttt{Init} is nondeterministic over \texttt{crasher} (one initial state per
subset of \texttt{SCENE}), and every one must be explored.

% ============================================================
\section{Fidelity: reproducing the loops}
\label{sec:fidelity}
% ============================================================

A model that cannot reproduce the known defect proves nothing.  Two corrective
clauses each guard a distinct loop, and each has a companion buggy variant that
strikes exactly that clause and reproduces the loop it prevents.

\paragraph{Quarantine (\texttt{display\_crash\_loop\_buggy.tex}).}
The variant strikes the quarantine effect from \texttt{BatchCrash} and
\texttt{IsolCrash} --- a death is tallied but no scene is ever quarantined
(\texttt{quarantined} stays empty).  Then the guard \texttt{s? $\notin$
quarantined} never excludes a crasher, \texttt{IsolCrash} fires forever against
the same scene, and \texttt{crashes} grows without bound.  This is the original
deterministic loop: render $\to$ crash $\to$ respawn $\to$ render.  ProB reaches
\texttt{crashes > 1 + (THRESHOLD-1)*card(crasher)}, the goal that is \emph{not
found} for the fixed design.

\paragraph{Stable-interval exit
(\texttt{display\_crash\_loop\_earlyexit\_buggy.tex}).}
The variant weakens \texttt{IsolExit}'s guard, dropping \texttt{crasher $\subseteq$
quarantined} so the worker may leave isolation \emph{early} --- the ``one clean
pass'' exit --- and clear the tallies of scenes still in play.  Now an intermittent
crasher is never caught: it crashes the batch (tallied), the worker enters
isolation, the crasher renders cleanly on its probe, \texttt{IsolExit} fires early
and wipes its tally, the worker returns to batching, and the next batched death
starts the cycle over.  \texttt{crashes} grows without bound.  This is the
intermittent-crasher loop the stable-interval exit exists to close, and its
reachability is the evidence that the \texttt{STABLE\_INTERVAL} $\geq$
\texttt{ATTRIBUTION\_WINDOW} tie is load-bearing, not decorative:
\begin{verbatim}
  # Both buggy variants (must report: goal FOUND)
  probcli docs/display_crash_loop_buggy.tex -model_check -nodead \
    -goal "crashes>1+(THRESHOLD-1)*card(crasher)" \
    -p DEFAULT_SETSIZE 2 -p MAX_INITIALISATIONS 20 -p MAXINT 8
  probcli docs/display_crash_loop_earlyexit_buggy.tex -model_check -nodead \
    -goal "crashes>1+(THRESHOLD-1)*card(crasher)" \
    -p DEFAULT_SETSIZE 2 -p MAX_INITIALISATIONS 20 -p MAXINT 8
\end{verbatim}

Each fixed-design goal that is \emph{not found} becomes \emph{found} under the
single struck clause, with the loop trace.  That difference --- the fixed design
terminates and each buggy design does not --- is the evidence that the model
detects the loops the two clauses were added to prevent.

\end{document}
